\chapter{Diseases of the Lungs and Respiratory System}
\label{ch:lungs}
The lungs are the primary organs of the respiratory system, responsible for gas exchange between the atmosphere and the bloodstream. The respiratory tract is divided into the upper airways (nose, pharynx, larynx) and lower airways (trachea, bronchi, bronchioles, and alveoli). Diseases of the respiratory system include obstructive and restrictive disorders, infections, neoplasms, and vascular conditions.

%-------------------------------------------------------------
\section{Acute Respiratory Infections}
%-------------------------------------------------------------

%-------------------------------------------------------------

%-------------------------------------------------------------

%-------------------------------------------------------------

%-------------------------------------------------------------

%-------------------------------------------------------------

\label{sec:Acute Respiratory Infections}
%-------------------------------------------------------------%-------------------------------------------------------------

\subsection{Acute Bronchitis}
\label{sec:Acute Bronchitis}
Acute bronchitis is an inflammation of the bronchial mucosa, most commonly caused by viral infections (rhinoviruses, coronaviruses, influenza, parainfluenza, adenovirus), accounting for approximately 90\% of cases; bacterial causes include \textit{Mycoplasma pneumoniae}, \textit{Chlamydophila pneumoniae}, and \textit{Bordetella pertussis}. The condition results from infection of the bronchial mucosa leading to inflammation and increased mucus production. Clinical features include acute onset of cough (productive or nonproductive), often following an upper respiratory infection, with or without sputum production, chest discomfort, and low-grade fever; the cough may persist for 2--3 weeks. Diagnosis is clinical, with chest X-ray typically normal (helping to differentiate from pneumonia). Management is symptomatic: rest, hydration, analgesics, cough suppressants for sleep-disturbing cough, and bronchodilators for wheezing; antibiotics are not indicated for most cases. Prognosis is excellent, with resolution in 1--3 weeks.

\subsection{COVID-19}
\label{sec:COVID-19}
Coronavirus disease 2019 (COVID-19) is an infectious disease caused by severe acute respiratory syndrome coronavirus 2 (SARS-CoV-2), first identified in Wuhan, China in December 2019 and declared a pandemic by the WHO in March 2020. SARS-CoV-2 is a betacoronavirus that uses the angiotensin-converting enzyme 2 receptor for cell entry. The condition results from viral infection causing a spectrum of illness from asymptomatic infection to severe pneumonia with acute respiratory distress syndrome, multi-organ failure, and death. Transmission occurs primarily through respiratory droplets and aerosols, with an incubation period of 2--14 days. Clinical features include fever, cough, fatigue, myalgia, anosmia or ageusia, headache, and dyspnea; severe disease features hypoxemia, bilateral pulmonary infiltrates, lymphopenia, elevated inflammatory markers, and cytokine storm; complications include ARDS, thromboembolism, myocarditis, acute kidney injury, neurological manifestations, and multisystem inflammatory syndrome in children. Diagnosis is by RT-PCR testing or rapid antigen tests. Management depends on severity: supportive care for mild disease; dexamethasone, remdesivir, baricitinib, tocilizumab, and anticoagulation for hospitalized patients; long COVID is characterized by persistent symptoms lasting weeks to months after acute infection. Prognosis depends on age, comorbidities, and disease severity; vaccination has significantly reduced severe disease and death.

\subsection{Influenza}
\label{sec:Influenza}
Influenza is an acute respiratory infection caused by influenza viruses (types A and B), belonging to the Orthomyxoviridae family; influenza A viruses are further classified by hemagglutinin and neuraminidase surface antigens, with antigenic drift causing seasonal epidemics and antigenic shift potentially causing pandemics. The condition results from viral infection of the respiratory epithelium. Transmission occurs via respiratory droplets and aerosols, with an incubation period of 1--4 days. Clinical features include abrupt onset of fever, chills, myalgia, headache, malaise, nonproductive cough, sore throat, and nasal congestion; complications include viral or secondary bacterial pneumonia, myocarditis, pericarditis, encephalopathy, myositis, and Reye's syndrome (in children given aspirin). Diagnosis is by rapid influenza diagnostic tests, molecular assays (RT-PCR), or viral culture. Management includes neuraminidase inhibitors (oseltamivir, zanamivir) and cap-dependent endonuclease inhibitor (baloxavir marboxil), most effective when initiated within 48 hours of symptom onset. Prognosis is excellent for most patients; annual influenza vaccination is the most effective preventive measure.

\subsection{Pleurisy}
\label{sec:Pleurisy}
Pleurisy is an inflammation of the pleural membranes causing sharp, stabbing chest pain that worsens with breathing, coughing, or sneezing; causes include viral infections (most common), bacterial pneumonia with parapneumonic effusion, pulmonary embolism, autoimmune diseases, malignancy, tuberculosis, and uremia. The condition results from inflammation of the parietal and visceral pleura. Clinical features include pleuritic chest pain (worse on inspiration, improved by holding breath) and possibly a pleural friction rub on auscultation; pleural effusion may develop. Diagnosis involves chest X-ray, thoracic ultrasound, diagnostic thoracentesis with Light's criteria, and CT chest for underlying pathology. Management targets the underlying cause; pleuritic pain is managed with NSAIDs; large symptomatic effusions may require therapeutic thoracentesis; recurrent pleurisy may necessitate pleurodesis. Prognosis depends on the underlying cause.

%-------------------------------------------------------------
\section{Lung Neoplasms}
%-------------------------------------------------------------

%-------------------------------------------------------------

Lung neoplasms represent the leading cause of cancer-related mortality worldwide, categorized into non-small cell lung cancer (85%) and small cell lung cancer (15%). Inhaled carcinogens, particularly cigarette smoke, induce cumulative genomic mutations leading to uncontrolled bronchial or alveolar epithelial proliferation. Early clinical silent progression delays diagnosis; management relies on surgical resection, platinum-based chemotherapy, targeted molecular therapy, and immunotherapy.

%-------------------------------------------------------------

%-------------------------------------------------------------

%-------------------------------------------------------------

%-------------------------------------------------------------

\label{sec:Lung Neoplasms}
%-------------------------------------------------------------%-------------------------------------------------------------

\subsection{Lung Cancer}
\label{sec:Lung Cancer}
Lung cancer is the leading cause of cancer-related death worldwide, with approximately 1.8 million deaths annually; non-small cell lung cancer accounts for 85\% of cases and includes adenocarcinoma (most common subtype, often in non-smokers, with driver mutations in EGFR, ALK, ROS1, KRAS), squamous cell carcinoma (central, strongly smoking-related), and large cell carcinoma; small cell lung cancer accounts for 15\%, is strongly associated with smoking, and has aggressive behavior with early metastasis. The condition results from malignant transformation of bronchial epithelial cells driven by carcinogen exposure. Clinical features include cough, hemoptysis, dyspnea, chest pain, weight loss, and post-obstructive pneumonia. Diagnosis involves CT-guided biopsy, bronchoscopy, or mediastinoscopy, followed by molecular testing for driver mutations. Management of early-stage disease is surgical resection with adjuvant chemotherapy; advanced disease is treated with platinum-based chemotherapy, immunotherapy, and targeted therapy for driver mutations; small cell lung cancer is treated with chemoradiation and prophylactic cranial irradiation. Prognosis depends on stage, histology, and molecular profile at diagnosis.

\subsection{Mesothelioma}
\label{sec:Mesothelioma}
Malignant pleural mesothelioma is a rare, aggressive tumor arising from the mesothelial cells of the pleura, strongly associated with asbestos exposure (latency 20--40 years); it is classified into epithelioid (most common, better prognosis), sarcomatoid, and biphasic subtypes. The condition results from asbestos fibers causing chronic inflammation and malignant transformation of mesothelial cells. Clinical features include progressive dyspnea, chest pain, weight loss, and pleural effusion. Diagnosis requires CT scan showing circumferential pleural thickening and pleural biopsy. Staging is by the TNM system. Management is palliative in most cases: chemotherapy with pemetrexed and cisplatin is first-line; immunotherapy has shown benefit; surgery may be considered in early-stage disease. Prognosis is poor, with median survival of 12--18 months.

%-------------------------------------------------------------
\section{Obstructive Lung Diseases}
%-------------------------------------------------------------

%-------------------------------------------------------------

%-------------------------------------------------------------

%-------------------------------------------------------------

%-------------------------------------------------------------

%-------------------------------------------------------------

\label{sec:Obstructive Lung Diseases}
%-------------------------------------------------------------%-------------------------------------------------------------

\subsection{Asthma}
\label{sec:Asthma}
Asthma is a heterogeneous chronic inflammatory disease of the airways characterized by wheezing, breathlessness, chest tightness, and cough, particularly at night or in the early morning; triggers include allergens, respiratory infections, exercise, cold air, and irritants. The condition results from IgE-mediated mast cell degranulation, eosinophilic infiltration, goblet cell hyperplasia, smooth muscle hypertrophy, and airway remodeling, leading to variable airflow obstruction and airway hyperresponsiveness. Clinical features include episodic wheezing, breathlessness, chest tightness, and cough. Diagnosis is supported by variable airflow limitation on spirometry (reversibility of more than 12\% and 200 mL in FEV$_1$ after bronchodilator), peak flow variability, and bronchoprovocation testing. Management follows stepwise therapy: inhaled corticosteroids (cornerstone of controller therapy), long-acting beta-agonists, leukotriene receptor antagonists, long-acting muscarinic antagonists, and biologic therapies (anti-IgE, anti-IL5, anti-IL4R$\alpha$) for severe eosinophilic asthma. Prognosis is excellent with appropriate controller therapy.

\subsection{Asthma-COPD Overlap}
\label{sec:Asthma-COPD Overlap}
Asthma-COPD overlap is a condition in which patients share features of both asthma and COPD, including persistent airflow limitation, airway inflammation that is both eosinophilic and neutrophilic, and significant exacerbation burden. The condition results from overlapping pathophysiological mechanisms of both asthma and COPD, with greater morbidity than either disease alone. Clinical features include persistent airflow limitation with variability and significant exacerbation burden. Diagnosis requires a history consistent with asthma (early onset, atopy, variability) combined with risk factors for COPD (smoking, occupational exposures). Management combines inhaled corticosteroids with long-acting bronchodilators; biologic therapies are considered for severe disease. Prognosis is worse than for asthma or COPD alone.

\subsection{Bronchiectasis}
\label{sec:Bronchiectasis}
Bronchiectasis is a permanent, abnormal dilation of the bronchi and bronchioles caused by chronic inflammation and recurrent infection; etiologies include cystic fibrosis, primary ciliary dyskinesia, immunodeficiency, prior severe respiratory infections, and idiopathic causes. The condition results from impaired mucociliary clearance, mucus stasis, chronic bacterial infection, and progressive airway destruction. Clinical features include chronic productive cough with copious purulent sputum, recurrent respiratory infections, hemoptysis, and progressive dyspnea. Diagnosis is by high-resolution CT showing tram-track opacities, signet ring sign, and dilated bronchi. Management includes airway clearance techniques, targeted antibiotics for exacerbations and chronic suppression, bronchodilators, and management of underlying causes; macrolide prophylaxis reduces exacerbation frequency. Prognosis depends on the underlying cause and frequency of exacerbations.

\subsection{Chronic Obstructive Pulmonary Disease}
\label{sec:Chronic Obstructive Pulmonary Disease}
Chronic obstructive pulmonary disease is a common, preventable, and treatable disease characterized by persistent respiratory symptoms and airflow limitation due to airway or alveolar abnormalities, usually caused by significant exposure to noxious particles or gases; cigarette smoking is the predominant risk factor (approximately 80\% of cases). The condition results from chronic inflammation leading to airway remodeling and alveolar destruction; COPD encompasses chronic bronchitis (persistent cough and sputum production) and emphysema (destruction of alveolar walls). Clinical features include progressive dyspnea, chronic cough, sputum production, wheezing, and chest tightness. Diagnosis is by spirometry showing post-bronchodilator FEV$_1$/FVC ratio less than 0.70; GOLD classification stages severity by FEV$_1$. Management includes smoking cessation (most important intervention), bronchodilators, inhaled corticosteroids for frequent exacerbators, pulmonary rehabilitation, long-term oxygen therapy for chronic hypoxemic failure, and lung volume reduction surgery or lung transplantation for advanced disease. Prognosis depends on disease severity, exacerbation frequency, and smoking cessation.

\subsection{Cystic Fibrosis}
\label{sec:Cystic Fibrosis}
Cystic fibrosis is an autosomal recessive disorder caused by mutations in the CFTR gene on chromosome 7, affecting approximately 1 in 2,500 Caucasian individuals. The condition results from a defective chloride channel leading to thick, viscous secretions in multiple organs. Pulmonary manifestations are the primary cause of morbidity and mortality, including chronic obstructive lung disease with bronchiectasis, chronic \textit{Pseudomonas aeruginosa} colonization, and progressive respiratory failure; pancreatic insufficiency causes malabsorption and failure to thrive; other complications include CF-related diabetes, liver disease, sinusitis, and male infertility. Clinical features include chronic cough, recurrent infections, and growth failure. Diagnosis is by newborn screening, sweat chloride testing (greater than 60 mmol/L), and CFTR genotyping. Management includes CFTR modulators that target the underlying protein defect, airway clearance, inhaled mucolytics, and antimicrobials. Prognosis has improved dramatically, with median survival of approximately 50 years in developed countries.

%-------------------------------------------------------------
\section{Pulmonary Vascular Diseases}
%-------------------------------------------------------------

%-------------------------------------------------------------

%-------------------------------------------------------------

%-------------------------------------------------------------

%-------------------------------------------------------------

%-------------------------------------------------------------

\label{sec:Pulmonary Vascular Diseases}
%-------------------------------------------------------------%-------------------------------------------------------------

\subsection{Pulmonary Embolism}
\label{sec:Pulmonary Embolism}
Pulmonary embolism is a potentially life-threatening condition caused by obstruction of the pulmonary arteries, most commonly by deep vein thrombosis of the lower extremities; risk factors include immobility, surgery, malignancy, pregnancy, oral contraceptives, obesity, and inherited thrombophilia. The condition results from thrombotic occlusion of pulmonary arteries. Clinical features range from dyspnea and pleuritic chest pain to hemoptysis, syncope, and hemodynamic instability (massive PE). Diagnosis involves clinical probability assessment (Wells score), D-dimer, and CT pulmonary angiography. Management depends on severity: anticoagulation for hemodynamically stable disease; systemic thrombolysis, catheter-directed therapy, or surgical embolectomy for massive disease with hemodynamic instability; inferior vena cava filters are placed when anticoagulation is contraindicated. Prognosis is excellent with prompt diagnosis and appropriate treatment.

\subsection{Pulmonary Hypertension}
\label{sec:Pulmonary Hypertension}
Pulmonary hypertension is defined as a mean pulmonary artery pressure of 20 mmHg or greater at rest, measured by right heart catheterization; the 2022 classification groups pulmonary hypertension into five categories: Group 1 (pulmonary arterial hypertension), Group 2 (left heart disease), Group 3 (lung disease or hypoxia), Group 4 (chronic thromboembolic disease), and Group 5 (unclear or multifactorial). The condition results from vascular remodeling with intimal fibrosis, medial hypertrophy, and plexiform lesions. Clinical features include progressive dyspnea, fatigue, syncope, and right heart failure. Diagnosis involves echocardiography, right heart catheterization, and evaluation of the underlying cause. Management of pulmonary arterial hypertension includes phosphodiesterase-5 inhibitors, endothelin receptor antagonists, prostacyclin analogs, and soluble guanylate cyclase stimulators; Group 4 disease is treated with pulmonary endarterectomy or balloon pulmonary angioplasty. Prognosis depends on the underlying group and response to therapy.

%-------------------------------------------------------------
\section{Respiratory Infections}
%-------------------------------------------------------------

%-------------------------------------------------------------

%-------------------------------------------------------------

%-------------------------------------------------------------

%-------------------------------------------------------------

%-------------------------------------------------------------

\label{sec:Respiratory Infections}
%-------------------------------------------------------------%-------------------------------------------------------------

\subsection{Community-Acquired Pneumonia}
\label{sec:Community-Acquired Pneumonia}
Community-acquired pneumonia is an acute infection of the lung parenchyma acquired outside of a hospital setting; common pathogens include \textit{Streptococcus pneumoniae} (most common typical bacterial cause), \textit{Haemophilus influenzae}, atypical organisms (\textit{Mycoplasma pneumoniae}, \textit{Chlamydophila pneumoniae}, \textit{Legionella pneumophila}), and respiratory viruses. The condition results from microbial infection of the alveolar spaces leading to consolidation. Clinical features include fever, productive cough, dyspnea, pleuritic chest pain, and malaise. Diagnosis is based on chest X-ray showing consolidation; the CURB-65 score assesses severity and guides disposition. Management includes empiric antibiotics: amoxicillin or macrolides for outpatient mild disease; respiratory fluoroquinolones or beta-lactam plus macrolide for inpatient disease; severe disease requires ICU admission and combination therapy. Prognosis depends on severity and patient factors.

\subsection{Empyema}
\label{sec:Empyema}
Empyema is the accumulation of purulent fluid within the pleural space, most commonly as a complication of pneumonia, thoracic surgery, or trauma; it progresses through exudative, fibrinopurulent, and organizational stages. The condition results from bacterial infection spreading to the pleural space. Clinical features include persistent fever despite appropriate antibiotic therapy for pneumonia, pleuritic chest pain, and dyspnea. Diagnosis is by chest X-ray, ultrasound, and diagnostic thoracentesis showing purulent fluid with low pH, low glucose, high LDH, and positive Gram stain or culture. Management includes chest tube drainage, intrapleural fibrinolytics (alteplase and DNase), VATS decortication for advanced disease, and appropriate antibiotics. Prognosis is good with prompt drainage and appropriate antibiotics.

\subsection{Lung Abscess}
\label{sec:Lung Abscess}
A lung abscess is a necrotizing infection of the lung parenchyma leading to a cavitary lesion containing purulent material; it may be primary (aspiration pneumonia in patients with impaired consciousness, poor dentition, or swallowing disorders) or secondary (bronchial obstruction, hematogenous seeding). The condition results from microbial infection leading to lung tissue necrosis. Common organisms include anaerobic bacteria, \textit{Staphylococcus aureus}, and \textit{Klebsiella pneumoniae}. Clinical features include fever, productive cough with foul-smelling sputum, weight loss, and night sweats. Diagnosis is by chest X-ray showing a thick-walled cavity, often with an air-fluid level. Management requires prolonged intravenous antibiotics (6--8 weeks), bronchoscopic drainage for large or refractory abscesses, and surgical resection as a last resort. Prognosis is good with appropriate antibiotic therapy.

\subsection{Tuberculosis}
\label{sec:Tuberculosis}
Tuberculosis is a chronic infectious disease caused by \textit{Mycobacterium tuberculosis}, transmitted via airborne droplets, and remains one of the leading infectious causes of death worldwide (approximately 1.3 million deaths annually). The condition results from infection with \textit{M. tuberculosis}, which may remain latent or progress to active disease. Primary tuberculosis may be asymptomatic or present with lower lobe infiltrates and hilar lymphadenopathy; reactivation tuberculosis typically affects the upper lobes with cavitary disease, hemoptysis, night sweats, weight loss, and fever. Clinical features of active disease include cough, hemoptysis, fever, night sweats, and weight loss. Diagnosis of active disease requires sputum smear and culture for acid-fast bacilli, nucleic acid amplification testing, and chest imaging; latent infection is detected by tuberculin skin testing or interferon-gamma release assays. Management of drug-susceptible disease involves a 6-month regimen of isoniazid, rifampicin, pyrazinamide, and ethambutol for 2 months, followed by isoniazid and rifampicin for 4 months; multidrug-resistant disease requires prolonged complex regimens. Prognosis is excellent for drug-susceptible disease with appropriate treatment.

%-------------------------------------------------------------
\section{Restrictive Lung Diseases}
%-------------------------------------------------------------

%-------------------------------------------------------------

%-------------------------------------------------------------

%-------------------------------------------------------------

%-------------------------------------------------------------

%-------------------------------------------------------------

\label{sec:Restrictive Lung Diseases}
%-------------------------------------------------------------%-------------------------------------------------------------

\subsection{Asbestosis}
\label{sec:Asbestosis}
Asbestosis is a chronic interstitial lung disease caused by the inhalation of asbestos fibers, typically occurring after occupational exposure (shipbuilding, construction, insulation manufacturing) with a latency period of 20--30 years. The condition results from macrophage phagocytosis of inhaled asbestos fibers, leading to chronic inflammation, fibrogenic growth factor release, and diffuse pulmonary fibrosis. Clinical features include progressive exertional dyspnea, dry cough, end-inspiratory bibasilar crackles, and digital clubbing. Diagnosis is based on occupational history, high-resolution CT showing pleural plaques, lower lobe predominant reticular opacities, and asbestos bodies (ferruginous bodies) on bronchoalveolar lavage or lung biopsy. Management is supportive, including smoking cessation, oxygen therapy, and vaccination; patients are at high risk for bronchogenic carcinoma and malignant mesothelioma. Prognosis is progressive and irreversible.

\subsection{Atelectasis}
\label{sec:Atelectasis}
Atelectasis is the collapse or incomplete expansion of pulmonary parenchyma, leading to impaired gas exchange. The condition results from airway obstruction by mucus plugs, foreign bodies, or tumors (resorptive atelectasis); pleural effusion or pneumothorax (compression atelectasis); or loss of pulmonary surfactant (microatelectasis). It is the most common cause of early postoperative fever (within 24--48 hours). Clinical features include tachypnea, cough, hypoxemia, and decreased breath sounds over the affected region. Diagnosis is by chest radiograph showing opacification with volume loss and tracheal deviation toward the affected side. Management involves incentive spirometry, chest physiotherapy, early postoperative ambulation, and therapeutic bronchoscopy for persistent mucus plugging. Prognosis is excellent when treated promptly.

\subsection{Idiopathic Pulmonary Fibrosis}
\label{sec:Idiopathic Pulmonary Fibrosis}
Idiopathic pulmonary fibrosis is a chronic, progressive, fibrosing interstitial pneumonia of unknown cause, occurring predominantly in older males. The pathology is characterized by usual interstitial pneumonia pattern: patchy, subpleural, basal-predominant fibrosis with honeycombing and fibroblastic foci. The condition results from aberrant wound healing and fibroblast proliferation in the lung parenchyma. Clinical features include insidious-onset exertional dyspnea and dry cough; Velcro-like crackles are heard on auscultation. Diagnosis requires HRCT showing UIP pattern and exclusion of known causes; lung biopsy may be needed when HRCT is inconclusive; pulmonary function tests show a restrictive pattern with reduced DLCO. Management includes antifibrotic therapies (pirfenidone and nintedanib) that slow disease progression; lung transplantation is the only intervention shown to improve survival. Prognosis is poor, with median survival of 3--5 years.

\subsection{Pneumoconiosis}
\label{sec:Pneumoconiosis}
Pneumoconiosis is a group of interstitial lung diseases caused by inhalation of mineral dusts and fibrogenic particles. Asbestosis results from asbestos fiber inhalation, causing basal-predominant interstitial fibrosis and pleural plaques, with a latency period of 10--20 years; silicosis (from silica dust) presents in chronic, accelerated, and acute forms; coal workers' pneumoconiosis results from coal dust exposure; berylliosis is a granulomatous disease from beryllium exposure. The condition results from accumulation of inhaled mineral dusts in the lungs, triggering chronic inflammation and fibrosis. Clinical features include progressive dyspnea and cough. Diagnosis is by history of occupational exposure, HRCT, and pulmonary function tests showing restriction with reduced DLCO. Management involves exposure removal and supportive care. Prognosis depends on the type and duration of exposure; these conditions are largely preventable through occupational safety measures.

\subsection{Pneumothorax}
\label{sec:Pneumothorax}

Pneumothorax is the accumulation of air in the pleural space, leading to partial or complete lung collapse. The condition results from rupture of subpleural apical blebs (primary spontaneous pneumothorax, typically in tall, thin young men), underlying lung disease (secondary spontaneous pneumothorax in COPD, cystic fibrosis), or chest trauma. Tension pneumothorax occurs when a one-way valve mechanism allows air to accumulate continuously under pressure, causing mediastinal shift and hemodynamic collapse. Clinical features include sudden-onset ipsilateral pleuritic chest pain and dyspnea, with hyperresonance and decreased breath sounds on the affected side. Diagnosis is by chest radiograph showing a visceral pleural line without peripheral lung markings; tension pneumothorax is a clinical diagnosis requiring immediate needle decompression at the second intercostal space in the midclavicular line prior to imaging. Management of small spontaneous pneumothoraces involves observation and oxygen administration; larger pneumothoraces require tube thoracostomy (chest tube insertion). Prognosis is good with prompt treatment.

\subsection{Sarcoidosis}
\label{sec:Sarcoidosis}
Sarcoidosis is a multisystem granulomatous disorder of unknown etiology characterized by non-caseating granulomas, most commonly affecting the lungs and intrathoracic lymph nodes; it disproportionately affects young adults, particularly African Americans and Scandinavians. The condition results from an aberrant immune response leading to granuloma formation. Pulmonary involvement (present in 90\% of cases) manifests as cough, dyspnea, and chest pain; extrapulmonary disease involves the skin (erythema nodosum, lupus pernio), eyes (uveitis), joints, heart, and nervous system. Clinical features depend on the organs involved. Diagnosis requires histological demonstration of non-caseating granulomas and exclusion of alternative causes; elevated ACE levels and hypercalcemia support the diagnosis; staging uses chest X-ray (Scadding stages). Management is indicated for symptomatic or progressive disease: corticosteroids are first-line, with methotrexate, azathioprine, or infliximab as steroid-sparing agents. Prognosis is favorable, with many cases resolving spontaneously.

\subsection{Silicosis}
\label{sec:Silicosis}
Silicosis is an occupational pneumoconiosis caused by chronic inhalation of crystalline silicon dioxide (silica) dust in miners, sandblasters, stonecutters, and foundry workers. The condition results from silica particle ingestion by alveolar macrophages, triggering cell death, release of fibrogenic cytokines, and formation of dense hyalinized silicotic nodules. Clinical features include exertional dyspnea, chronic productive cough, and increased susceptibility to pulmonary tuberculosis (silicotuberculosis) due to impaired macrophage function. Diagnosis is supported by chest imaging showing upper lobe eggshell calcification of hilar lymph nodes and nodular opacities. Management is preventive and supportive; oxygen therapy and prompt treatment of mycobacterial infections are essential. Prognosis depends on cumulative silica exposure.

